\chapter{Introduction}
\label{Chap1}

\section{Background}

Large vision--language models (LVLMs) process images together with
natural-language instructions. They can classify and describe images as well
as reason about depicted actions, relationships, and social situations
\cite{zellers2019recognition}. When an image has social or normative
significance, however, identifying its visible objects is only the first step.
The model must also determine what is happening and which moral concern, if
any, is supported by the scene. This makes visual moral classification
particularly challenging when the relevant context is incomplete or
ambiguous.

Two main challenges arise from this task. First, some of the information needed
for a moral judgement may not be visible. An image may show an action without
revealing the actor's intention, the relationship between the people involved,
the events preceding the scene, or its consequences. For example, physical
contact could represent assistance, restraint, aggression, or an accidental
interaction. These interpretations could lead to different moral categories,
even though the visible action is the same. Visual information nevertheless
contains evidence relevant to human moral judgement
\cite{zhu2025visual}. The challenge is therefore to use that evidence without
presenting an inferred context as something directly observed in the image.

Second, different models may resolve the same ambiguity differently. LVLMs are
developed using different training data, architectures, and post-training
objectives. They may therefore emphasise different visual cues, make different
assumptions about missing context, or apply different thresholds when assigning
a moral category. Previous comparisons have shown that moral-foundation
profiles and moral choices can vary across language models and elicitation
conditions \cite{abdulhai2024foundations,liu2024pluralistic}. The same image
and classification instructions may consequently produce different moral
judgements across models.

A common category framework is required to compare these judgements. Moral
Foundations Theory (MFT) distinguishes five paired dimensions of moral concern:
Care/Harm, Fairness/Cheating, Loyalty/Betrayal, Authority/Subversion, and
Sanctity/Degradation
\cite{graham2009liberals,graham2011mapping}. These dimensions provide a
consistent vocabulary for describing the type of moral concern a model assigns
to an image. In this study, the five pairs are represented by ten selectable
categories, together with a \textit{None of these / Not clearly moral in
nature} option. MFT is used to make model outputs comparable, not to determine
a universally correct label. This distinction is important because human
moral judgements can vary across cultural and gender groups
\cite{qian2024cultural}.

Recent research has begun to evaluate moral judgement in multimodal systems
using several different task designs. MM-MoralBench uses synthetic
image--dialogue scenarios to evaluate moral judgement and classification
\cite{yan2026mmmoralbench}. MORALISE instead examines moral judgement and norm
attribution in curated real-world image--text data
\cite{lin2026moralise}. Other studies compare moral narratives produced by
people and models from similar visual information
\cite{rezapour2025tales}, or elicit graded judgements grounded in multimodal
evidence \cite{park2026mmscale}. Taken together, this work shows that
multimodal moral evaluation is feasible and provides the foundation for
further comparison of how LVLMs classify morally relevant images.

\section{Problem Statement and Research Gap}

Despite recent progress, evaluating LVLMs in visual moral classification
remains challenging. Existing evaluations often summarise performance using
aggregate measures such as accuracy or agreement. Although these measures are
useful for comparison, they do not show how models behave on individual
images. In morally ambiguous cases, two models may achieve similar overall
scores even when they assign different categories to the same images or
respond differently to changes in visible content.

Previous studies examine multimodal moral judgement using different datasets,
tasks, taxonomies, and reference judgements
\cite{yan2026mmmoralbench,lin2026moralise,rezapour2025tales,park2026mmscale}.
Consequently, it remains unclear how multiple deployed LVLMs compare when they
classify the same images using the same moral categories. It is also unclear
how their classifications correspond to the range of participant judgements
and whether their outputs differ after selected visual information is
modified. The research gap is therefore the absence of an integrated
comparison of model classifications, participant responses, and responses to
selected visual modifications within a common evaluation setting.

To address this gap, this dissertation compares three LVLM configurations on
a shared image corpus using the same moral-category framework. It examines
their overall classification patterns, identifies where their classifications agree or
differ, compares their classifications with participant responses, and
examines their outputs in selected edited-image cases. The repeatability
component of RQ2 also checks whether classifications remain stable across repeated
runs. Together, these analyses reveal differences that a single aggregate
agreement score would not capture.

\section{Research Aim and Questions}

The aim of this dissertation is to compare how three LVLM configurations
classify images using a shared moral-category framework. Specifically, the
study examines their overall classification patterns, their agreement and
disagreement, the correspondence between model and participant judgements,
and variation in model outputs across selected original and edited image
pairs. This aim is addressed through four research questions:

\begin{description}

\item[\textbf{RQ1:}] What classification patterns do GPT-4o, GPT-5.6, and
Qwen3.6 Flash produce within the fixed 127-image corpus under the visual
moral-category classification task?

\item[\textbf{RQ2:}] To what extent do the three models agree with one
another, and how stable are their classifications across repeated runs?

\item[\textbf{RQ3:}] How do participants differ in their moral-category
selections, and to what extent do their selections correspond to those of
the three model configurations across the survey images?

\item[\textbf{RQ4:}] What changes in moral-category labels and model-reported
confidence are observed between selected base images and their edited
variants?

\end{description}

\section{Scope and Contributions}

This dissertation examines the observable classification behaviour of three
LVLM configurations: GPT-4o, GPT-5.6, and Qwen3.6 Flash. Each configuration
applies the same eleven-category moral framework to a fixed corpus of 127
images and returns a category label, a model-reported confidence score, and a
brief reason. Because the corpus has no independently validated image-level
ground truth, the study describes classification patterns and agreement rather
than measuring accuracy or determining whether one model is morally correct.

Three additional analyses extend the primary model comparison. First, a
participant study uses 56 completed responses from a fixed 16-image pool to
compare model-selected categories with participants' multi-label selections.
Second, an exploratory image-edit analysis compares eleven modified variants
with their eight base images. Third, the stability component of RQ2 examines
repeated classifications on a stratified 32-image subset. The additional
Part~2 survey ratings are reported separately and do not answer the four main
research questions.

The study makes three main contributions. First, it provides a systematic
comparison of overall classification patterns, inter-model agreement, and
output stability under a shared task. The results show that similar aggregate
rates within this corpus do not necessarily imply image-level agreement on
individual moral categories.
Second, it compares model classifications with participant responses while
preserving variation in participants' multi-label judgements rather than
treating them as a single correct answer. This distinguishes agreement between
models from correspondence between models and people. Third, it uses selected
original and edited image pairs to identify cases in which changes to visible
content coincide with changes in model labels or model-reported confidence.
This connects the aggregate comparison with case-level sensitivity to visual
information.

The scope of these contributions is limited in two main respects. First, the
fixed image corpus is modest and constructed, its current researcher-coded
indexing metadata are strongly uneven, and it is not a representative sample.
Second, the participant study uses a convenience
sample, while the edited-image and stability analyses use restricted subsets
of the corpus. The image edits also do not isolate individual visual features
under fully controlled conditions. The findings should therefore be
interpreted as descriptive evidence from the evaluated models, participants,
and images, rather than as population-level model rankings or causal effects
of particular visual cues.


\section{Relevance to the Sustainable Development Goals}
\label{sec:sdg_relevance}

The study's relevance to the United Nations Sustainable Development Goals
(SDGs) is methodological. Its primary connection is to
SDG~16: Peace, Justice and Strong Institutions, particularly Target~16.6 on
accountable and transparent institutions and Target~16.7 on inclusive and
representative decision-making~\cite{un2015agenda}. If LVLM outputs are used to inform
institutional judgements, their treatment of differing moral interpretations
becomes relevant to accountability and inclusion.

SDG~10: Reduced Inequalities provides a secondary, indirect connection~\cite{un2015agenda}. Reducing
culturally or socially diverse moral interpretations to one label could
misrepresent particular people or contexts.

The dissertation does not assess an institutional deployment, measure unequal
outcomes, or test an SDG indicator. The connection therefore lies in how
interpretive variation is evaluated; the study does not measure progress
towards either goal.


\section{Dissertation Structure}

Chapter~2 reviews research on moral judgement, visual moral classification,
multimodal moral evaluation, model stability, cultural variation, and
subjective disagreement. Chapter~3 describes the image corpus, moral-category
framework, model inference procedure, participant survey, image-edit study,
and analytical methods. Chapter~4 presents and interprets the findings in the order of the four
research questions, followed by the supplementary prompt-robustness check.
Chapter~5 presents the conclusion, scope and limitations, directions for future
work, and implications for the Sustainable Development Goals.
